\documentclass[11pt]{article}
\usepackage[margin=1in]{geometry}
\usepackage{booktabs}
\usepackage{longtable}
\usepackage{array}
\usepackage{caption}
\usepackage{graphicx}
\usepackage[utf8]{inputenc}

\title{Experiment Results}
\author{}
\date{}

\begin{document}
\maketitle

\section{Single-AIS Trajectory Forecasting Comparison}

Table~\ref{tab:single-ais-main} summarizes the current test-set comparison on the traj\_a-only single-AIS setting.
The metrics \texttt{Avg L2}, \texttt{Avg MSE}, and \texttt{Avg DTW} are computed after reconstructing predictions back to latitude/longitude space.
Lower is better for all reported metrics.

\begin{table}[htbp]
    \centering
    \caption{Test-set comparison on the traj\_a-only single-AIS setting.}
    \label{tab:single-ais-main}
    \begin{tabular}{lccccccc}
        \toprule
        Model & Epochs & Test Loss & Pos MAE & Final MAE & Avg L2 & Avg MSE & Avg DTW \\
        \midrule
        Ours   & 50 & \textbf{0.085999} & 147.2423 & \textbf{233.5597} & \textbf{0.002117} & \textbf{0.000006039} & \textbf{0.020905} \\
        AGDN   & 20 & 0.104537 & 162.5801 & 294.1568 & 0.002403 & 0.000006808 & 0.024231 \\
        TLSTM  & 20 & 0.087601 & \textbf{145.7649} & 240.4829 & 0.002154 & 0.000006595 & 0.021198 \\
        TPatch & 20 & 0.088869 & 149.7027 & 242.8796 & 0.002210 & 0.000006288 & 0.021267 \\
        \bottomrule
    \end{tabular}
\end{table}

\paragraph{Current observation.}
Under the current setting, our model achieves the best overall performance on Test Loss, Final MAE, Avg L2, and Avg DTW.
TLSTM is competitive on Pos MAE, but remains worse on the trajectory-level and endpoint-oriented metrics.

\begin{figure}[htbp]
    \centering
    \includegraphics[width=\textwidth]{workdir_singleais_traj_a/experiments/patch_nomissing_traj_a_singleais_gpu2_lr1e4_e50_slide1_20260406_181745/visualizations/test_ten_short_samples.png}
    \caption{Visualization of 10 test samples. Each panel shows the last 6 observed input steps, together with 6-step ground-truth future trajectory and 6-step predicted trajectory.}
    \label{fig:single-ais-ten-samples}
\end{figure}

\section{Mask Robustness Evaluation}

In this experiment, the input history is randomly masked from 10\% to 90\% with a step size of 10\%.
For our model, each masked time step is replaced by a temporally nearby observation from \texttt{traj\_b}.
For AGDN, TLSTM, and TPatch, masked time steps are directly filled with zeros.
The reported \texttt{L2}, \texttt{MSE}, and \texttt{DTW} values are all computed in reconstructed latitude/longitude space.

\input{workdir_singleais_traj_a/experiments/patch_nomissing_traj_a_singleais_gpu2_lr1e4_e50_slide1_20260406_181745/mask_robustness/mask_robustness_results.tex}

\section{Planned Experiments}

The following experiments are reserved for future completion:

\begin{itemize}
    \item Ablation study on the \texttt{traj\_b}-based completion strategy
    \item Influence of multi-source AIS data on masked-trajectory forecasting performance
    \item Comparison under different masking patterns (random mask, block mask, tail mask)
    \item Effect of different training data sizes
    \item Stability across multiple random seeds
    \item Error analysis on hard trajectory cases and failure examples
\end{itemize}

\end{document}
